\chapter{Conclusion}

Ce travail visait à explorer le lien intrinsèque entre la structure convolutionnelle des réseaux 
de neurones (CNN) et l'émergence de directions de vulnérabilité privilégiées dans le domaine de 
Fourier. Plus précisément, notre démarche a consisté à déterminer dans quelle mesure les propriétés 
mathématiques inhérentes aux couches de convolution — au-delà des biais liés aux données ou au 
processus d'entraînement — permettent d'expliquer l'existence de perturbations adversariales universelles.

Pour ce faire, après avoir posé le cadre général des CNN et des attaques adversariales, nous 
avons reconstruit et analysé le raisonnement théorique proposé par Tsuzuku et 
Sato~\cite{tsuzuku2019structural}, en nous appuyant initialement sur un modèle linéaire simplifié.

Afin de transposer ces résultats théoriques à l'architecture réelle des CNN, où la 
présence d'activations ReLU et d'opérations de \textit{max-pooling} empêche toute 
diagonalisation globale, nous avons introduit le concept de linéarisation locale via 
l'analyse du jacobien du modèle. Cette extension théorique s'est avérée cruciale pour 
justifier l'application de résultats fondés sur la linéarité à des architectures 
fondamentalement non linéaires.

L'ensemble de ces observations permet ainsi d'élucider le mécanisme par lequel une 
perturbation fondée sur une unique fréquence de Fourier peut impacter si sévèrement 
un réseau convolutionnel. L'amplification successive de cette perturbation au travers 
des différentes couches offre une explication structurelle probante à l'existence des 
perturbations universelles.

Notre démarche expérimentale vient conforter cette intuition. Les \textit{heatmaps} 
générées sur l'architecture \texttt{VGGSmallBN}, pour les jeux de données CIFAR-10 et 
CIFAR-100, mettent en évidence une sensibilité fréquentielle fortement hétérogène. 
Certaines fréquences engendrent en effet des \textit{fool ratios} significativement 
supérieurs, démontrant que des perturbations d'amplitude égale n'ont pas un impact 
équivalent : leur direction dans l'espace des entrées est un facteur déterminant. 
La comparaison de ces résultats avec l'injection d'un bruit aléatoire de même amplitude 
corrobore pleinement cette asymétrie.

Par ailleurs, les résultats soulignent l'influence logique de la complexité de la tâche 
de classification sur la robustesse du modèle. Le réseau s'avère globalement plus 
vulnérable sur CIFAR-100 que sur CIFAR-10. Cette fragilité, examinée plus en détail 
dans l'annexe~\ref{annexe:marges-logits}, est cohérente avec l'exigence accrue du problème : 
face à un nombre de classes supérieur, les frontières de décision se multiplient et se 
resserrent en raison de la proximité visuelle des catégories, permettant à une perturbation 
de modifier plus aisément la prédiction. Néanmoins, en dépit de cette vulnérabilité accrue, 
la sensibilité conserve une topologie structurée où certaines directions de Fourier demeurent prépondérantes.

Il faut toutefois rester prudent sur la portée de ces résultats. Notre reproduction expérimentale 
se limite à une architecture de type VGG et à des images de faible résolution. De plus, les 
fréquences les plus sensibles ont été choisies à partir de \textit{heatmaps} calculées sur 
l'ensemble de test. On ne peut donc pas affirmer que tous les réseaux convolutionnels présentent 
exactement la même structure de vulnérabilité. En revanche, dans le cadre étudié ici, les directions 
de Fourier apparaissent comme un outil pertinent pour mieux comprendre certaines sensibilités des CNN.

Cette analyse peut aussi être reliée à la question de l'explicabilité des modèles. En effet, lorsqu'une 
fréquence particulière modifie fortement les prédictions du réseau, cela indique que cette composante 
de l'image joue un rôle important dans sa décision. Autrement dit, les fréquences sensibles ne 
servent pas seulement à construire des perturbations adversariales : elles donnent aussi des indices 
sur le type d'information que le réseau utilise pour classifier les images.

Une suite possible à ce travail serait donc d'utiliser ces fréquences sensibles à la fois pour renforcer 
les modèles et pour mieux les interpréter. Par exemple, elles pourraient servir à guider de nouvelles 
formes d'augmentation de données, à limiter certaines amplifications internes, ou à comparer des 
architectures plus robustes, notamment avec des mécanismes d'anti-aliasing.

Cette direction permettrait de faire le lien entre trois aspects importants : la robustesse adversariale, 
la résistance aux variations des données, et l'interprétation des décisions prises par le réseau. 
L'objectif ne serait donc pas seulement de montrer que certaines fréquences de Fourier sont dangereuses 
pour le modèle, mais aussi de comprendre comment leur identification peut aider à construire des réseaux 
plus stables, plus robustes et plus compréhensibles.

